\documentclass[11pt]{article}
\usepackage{../../estilo/vanguarda44}

\renewcommand{\contentsname}{Sumário}

\title{\bfseries\color{vinho}\Huge VANGUARDA 44: FRENTE LESTE}
\author{}
\date{}

\hypersetup{
  pdftitle={Vanguarda 44: Frente Leste - Mini-Livro de Expansão},
  pdfauthor={Felipe Ramires Terrazas},
}

\fancyhead[L]{\small\color{cinzachumbo}Vanguarda 44: Frente Leste (Mini-Livro de Expansão)}

\begin{document}

\begin{titlepage}
\centering
\vspace*{3cm}
{\Huge\bfseries\color{vinho} VANGUARDA 44: FRENTE LESTE \par}
\vspace{0.5cm}
{\Large\color{mostarda} Mini-Livro de Expansão \par}
\vspace{2cm}
{\large Doutrina Soviética e o Cenário de Stalingrado \par}
\vfill
\end{titlepage}

\tableofcontents

\bigskip

A Frente Leste foi a mais brutal e decisiva da Segunda Guerra Mundial: um teatro de guerra de escala e ferocidade que a Frente Ocidental nunca chegou a igualar. Este mini-livro traz a doutrina soviética e o cenário mais simbólico da guerra: Stalingrado.

\section{Doutrina Nacional: União Soviética}

O Exército Vermelho compensava a falta de treinamento individual com número, artilharia massiva e ferocidade no combate corpo a corpo: o icônico grito de guerra ``Urrah!'' no assalto à baioneta.

\textbf{Regra do Exército}: esquadrões soviéticos ganham +1 dado extra em Combate Corpo a Corpo (CQC).

\subsection{Sub-Facção A: Exército Vermelho Regular}

Foco: estrutura de comando duplicada. Regra Especial (``Comissário Político''): o esquadrão pode incluir um Comissário Político como uma segunda fonte fixa de 1 PC, independente do Sargento. Mesmo que o Sargento morra, o esquadrão nunca fica sem comando enquanto o Comissário sobreviver.

\subsection{Sub-Facção B: Guardas Soviéticos (Elite)}

Foco: veteranos endurecidos em combate urbano. Regra Especial (``Veteranos de Stalingrado''): podem rolar novamente 1 dado de Dano perdido por ativação de disparo.

\section{Cenário: Stalingrado}

\textbf{Contexto Histórico}: A batalha por Stalingrado se tornou sinônimo de combate urbano implacável, travado prédio por prédio, andar por andar, entre o 6º Exército alemão e as forças soviéticas defendendo a cidade às margens do rio Volga: um dos pontos de virada mais decisivos de toda a guerra.

\textbf{Data e Local}: agosto de 1942 a fevereiro de 1943, Stalingrado, União Soviética.

\textbf{Forças}: Alemães (Atacante) vs. Soviéticos (Defensor).

\textbf{Terreno}: urbano extremo, praticamente toda a mesa em Cobertura Pesada ou dentro de edifícios, com ruínas por toda parte.

\textbf{Implantação}: Alemães avançam de uma borda; Soviéticos defendem em profundidade, podendo começar com unidades já em Ocupação Parcial de edifícios (regra da Parte V do Livro principal).

\textbf{Qualidade de Tropa}: Alemães Veteranos; Soviéticos Regulares (compensam com o bônus de CQC da doutrina nacional).

\textbf{Regra Especial ``Não Há Terra Além do Volga''}: unidades soviéticas nunca recuam voluntariamente: a ordem Correr não pode ser usada para se afastar do lado alemão da mesa (só para flanquear ou avançar).

\textbf{Objetivo}: Soviéticos vencem se ainda controlarem ao menos 1 de 2 a 3 edifícios-chave definidos na mesa até o turno 10. Alemães vencem se tomarem todos.

\textbf{Duração}: 10 turnos.

\end{document}
